\chapter{M007 --- Interface Prototype Construction}

\section{Stage 1 --- Revised Mathematical Specification}

\subsection{Scientific Purpose}

Construct an admissible interface prototype from two realized periodic slabs by
solving the crystallographic interface matching problem. This specification
supersedes the earlier candidate ontology in light of subsequent repository
analysis.

\subsection{Scientific Transformation}

\[
\mathrm{PeriodicSlab}_A \times \mathrm{PeriodicSlab}_B
\longrightarrow
\mathrm{InterfacePrototype}
\]

This is the first binary scientific transformation in the CALM workflow.

\subsection{Scientific Inputs}

Two realized periodic slabs.

\subsection{Scientific Output}

An Interface Prototype representing a crystallographically admissible candidate
geometry relating the two periodic slabs. The prototype is not yet a fully
constructed interface.

\subsection{Mathematical Interpretation}

An Interface Prototype records an admissible correspondence between two
periodic slabs after lattice compatibility constraints have been satisfied.
Ranking, Pareto filtering, grouping, canonicalization, and hashing are
subsequent operations on interface prototypes and are not part of this
scientific transformation.

\subsection{Repository Evidence}

Repository inspection indicates:

\begin{itemize}
\item Tutorial 3 defines an interface prototype as a candidate geometry with
commensurate lattices.
\item Prototype search consists of supercell enumeration, lattice matching,
strain analysis, and ranking.
\item \texttt{calm/interface/prototypes.py} performs canonicalization,
hashing, and deduplication of existing prototypes rather than constructing
them.
\item Prototype grouping operates on previously generated prototype results.
\end{itemize}

\subsection{Primary Audit Question}

Does the implementation faithfully construct Interface Prototype objects from
pairs of Periodic Slab objects while keeping representation and workflow
operations separate from the scientific transformation?
